\section{Discussion}
\label{sec:discussion}

\subsection{Why the Two Objectives Are Not Symmetric}

The group term moves its target by roughly $80\%$; the individual term moves
its own by roughly $5\%$. The asymmetry is not a tuning artefact but follows
from what each objective can control.

\softdp{} in \eqref{eq:softdp} is a difference of group means of $p$. A model
can reduce it by shifting the decision surface between groups---an operation
that costs little accuracy because it redistributes predictions rather than
requiring better ones. The quantity is, in effect, freely adjustable.

\softge{} in \eqref{eq:softge} is the inequality of $b_i = 1 + p_i - y_i$.
Benefit equality is attained exactly when $b_i$ is constant, that is when
every individual is predicted equally correctly. Its floor is therefore set by
the irreducible error of the hypothesis class on the data: one cannot equalize
benefit without in fact being more accurate. On COMPAS, where the ceiling is
approximately $0.68$ accuracy, there is little room, and the term accordingly
achieves little. A useful implication is that individual fairness in this
outcome-based sense is not primarily an allocation problem but a
\emph{prediction-quality} problem.

\subsection{$\beta$ as a Policy Parameter, Not a Hyperparameter}

Hyperparameters are tuned: one searches for the value that maximises a
criterion. $\beta$ does not admit that treatment, because the two ends of its
range optimise different criteria and there is no metric that ranks them.
Selecting $\beta$ by validation performance simply encodes whichever fairness
notion the validation metric happens to measure---which is exactly what our
utopia-distance rule does, and why it is scale-sensitive
(Section~\ref{sec:german}).

$\beta$ is better understood as a \emph{policy dial}: it makes explicit a
commitment that is otherwise implicit and undocumented in a deployed system.
Our results give that commitment a price. On Adult, moving from $\beta=0$ to
$\beta=1$ buys a $0.18$ reduction in demographic parity difference at a cost
of $0.22$ in equalized odds difference and $0.015$ in benefit inequality. An
institution that must state which definition it is optimizing can read the
consequences of that choice off the trade-off map rather than discovering them
after deployment.

\subsection{The Uncomfortable Case}

The Adult MLP at $\beta=1$ makes the tension concrete. It achieves
$\dpdtwo = 0.009$: men and women are selected at nearly identical rates, and
by the most widely cited group-fairness definition the model is close to
ideal. The same model has $\mathrm{EOD} = 0.321$, meaning its error rates
differ across sexes by roughly a third---worse than the unconstrained baseline
by a factor of $3.3$. Both statements are true of the same predictor.

There is no analysis internal to the data that resolves this. Whether it is
better to select equally or to err equally is a normative question about what
the decision is for. The value of a trade-off map is not that it answers the
question but that it prevents the answer from being made accidentally, by a
default in a fairness library.

\subsection{Where the Composite Loss Sits Among Existing Methods}

On group-fairness metrics the loss matches rather than beats
Reweighing~\cite{kamiran2012data},
ExponentiatedGradient~\cite{agarwal2018reductions} and
ThresholdOptimizer~\cite{hardt2016equality}
(Table~\ref{tab:master}, Fig.~\ref{fig:pareto}). We regard that as sufficient
for its purpose: an instrument used to measure a trade-off must be credible on
the axis it shares with existing tools, and it is.

Its distinguishing property is the one those methods lack. Each of the
baselines is constructed around a single constraint---demographic parity or
equalized odds---and offers no continuous path to an individual-fairness
objective. The composite loss traverses that path within a fixed architecture,
optimizer and evaluation protocol, so differences along $\beta$ are
attributable to the objective rather than to a change of method. That is what
makes Fig.~\ref{fig:beta} interpretable as a trade-off curve.

\subsection{A Sample-Size Floor for Fairness Model Selection}

The German Credit analysis (Section~\ref{sec:german}) generalizes beyond this
paper. Fairness metrics are differences of \emph{conditional} rates, so their
effective sample size is governed by the smallest group, not by $n$. With
$310$ women in the dataset and roughly $62$ in a validation split, the
standard error of a group selection rate is comparable to the entire disparity
being measured, and validation-based tuning ranks noise
(Table~\ref{tab:selection}).

Practitioners frequently select fairness hyperparameters exactly this way. Our
results suggest that such a procedure should be accompanied by a check---the
validation--test correlation of the fairness metric across configurations is a
cheap one---and that where it fails, a pre-specified configuration is
preferable to a tuned one. The cost of ignoring this is not merely a weaker
result: on German the tuned model reproduced the baseline disparity while a
pre-specified $\beta=1$ reduced it by $48\%$ at no accuracy cost.
